Distinct dynamical mechanisms can predict similar actions, so family recovery, predictive accuracy and practical loss equivalence must be evaluated separately. We examine these distinctions using synthetic feedback trajectories. Historical ARC-inspired stress tests motivate the problem but are not confirmatory evidence. A standardized development benchmark crosses six two-level observation factors with eight generator families, producing 40,960 trajectories with separate fitting, validation choice and final evaluation. A further 5,760 targeted trajectories include exact reduction of PID to P. In the full grid's reference arm, nominal PID recovery was 79/80; analyst-noise SD 0.04 reduced it to 41/80. High family recovery could coexist with predictive loss differences inside the development tolerance. Short and sparse observation schedules also left 2,432 attempts without enough observations, which remain in recovery denominators. A further exploratory extension generated 26,880 trajectories for independent-role validation, AIC/BIC comparison and observation stress tests. Penalization improved recovery of simpler nested labels but could reduce noisy PID recovery; bootstrap conclusions remained strongly tolerance-dependent. These development results establish conditional design sensitivity within a specified generator universe, not structural equivalence or human mechanisms. Final inferential claims require calibrated tolerance, precision planning and independent reproduction of a frozen numerical protocol.
